\begin{thebibliography}{99}
\small

\bibitem{who} World Health Organization. Cardiovascular Diseases. Webpage, 2023. Available at: \url{https://www.who.int/health-topics/cardiovascular-diseases} (Accessed on June 23, 2023).

\bibitem{mathers} Colin D. Mathers and Dejan Loncar. Projections of global mortality and burden of disease from 2002 to 2030. \emph{PLoS Med}, 3(11):e442, 2006. doi: \href{https://doi.org/10.1371/journal.pmed.0030442}{10.1371/journal.pmed.0030442}.

\bibitem{mayo} Mayo Clinic. Heart disease. Webpage, 2019. Available at: \url{https://www.mayoclinic.org/diseases-conditions/heart-disease/symptoms-causes/syc-20353118} (Accessed on June 23, 2023).

\bibitem{libby} Peter Libby. Mechanisms of acute coronary syndromes and their implications for therapy. \emph{New England Journal of Medicine}, 368(21):2004--2013, 2013. doi: \href{https://doi.org/10.1056/NEJMra1216063}{10.1056/NEJMra1216063}.

\bibitem{yusuf} Salim Yusuf, Steven Hawken, Stephanie \^{O}unpuu, Tony Dans, Alvaro Avezum, Fernando Lanas, Matthew McQueen, Andrzej Budaj, Prem Pais, John Varigos, and Liu Lisheng. Effect of potentially modifiable risk factors associated with myocardial infarction in 52 countries (the INTERHEART study): Case-control study. \emph{The Lancet}, 364(9438):937--952, 2004. doi: \href{https://doi.org/10.1016/S0140-6736(04)17018-9}{10.1016/S0140-6736(04)17018-9}.

\bibitem{virani} Salim S. Virani, Alvaro Alonso, Eunsik J. Benjamin, Marcio S. Bittencourt, Clifton W. Callaway, April P. Carson, Alanna M. Chamberlain, Alexander R. Chang, Susan Cheng, Francesca N. Delling, et al. Heart disease and stroke statistics---2020 update: A report from the American Heart Association. \emph{Circulation}, 141(9):e139--e596, 2020. doi: \href{https://doi.org/10.1161/CIR.0000000000000757}{10.1161/CIR.0000000000000757}.

\bibitem{alaref} Subhi J. Al'Aref, Kimberly Anchouche, Gaurav Singh, Patricia J. Slomka, Krzysztof K. Kolli, Amit Kumar, Mohit Pandey, Gabriel Maliakal, Alexander R. van Rosendael, Ashley N. Beecy, et al. Clinical applications of machine learning in cardiovascular disease and its relevance to cardiac imaging. \emph{European Heart Journal}, 40(24):1975--1986, 2019. doi: \href{https://doi.org/10.1093/eurheartj/ehy404}{10.1093/eurheartj/ehy404}.

\bibitem{gavhane} Aditi Gavhane, Gouthami Kokkula, Isha Pandya, and Kailas Devadkar. Prediction of heart disease using machine learning. In \emph{2018 Second International Conference on Electronics, Communication and Aerospace Technology (ICECA)}, pages 1275--1278, 2018. doi: \href{https://doi.org/10.1109/ICECA.2018.8474922}{10.1109/ICECA.2018.8474922}.

\bibitem{sheikabdullah} A. SheikAbdullah. A data mining model to predict and analyze the events related to coronary heart disease using decision trees with particle swarm optimization for feature selection. \emph{International Journal of Computer Applications}, 55:49--55, 2012. doi: \href{https://doi.org/10.5120/8779-2736}{10.5120/8779-2736}.

\bibitem{currie} Geoffrey Currie, Kristina Hawk, Eric Rohren, Alanna Vial, and Ran Klein. Machine learning and deep learning in medical imaging: Intelligent imaging. \emph{Journal of Medical Imaging and Radiation Sciences}, 50, 2019. doi: \href{https://doi.org/10.1016/j.jmir.2019.09.005}{10.1016/j.jmir.2019.09.005}.

\bibitem{chlif} M. Chlif, N. M. Ammar, N. Ben Said, L. Sergey, S. Ahmaidi, F. Alassery, and H. Hamam. Mechanism of dyspnea during exercise in children with corrected congenital heart disease. \emph{Int. J. Environ. Res. Public Health}, 19(99), 2022. doi: \href{https://doi.org/10.3390/ijerph19010099}{10.3390/ijerph19010099}.

\bibitem{hussain} N. M. Hussain, A. U. Rehman, M. T. B. Othman, J. Zafar, H. Zafar, and H. Hamam. Accessing artificial intelligence for fetus health status using hybrid deep learning algorithm (AlexNet-SVM) on cardiotocographic data. \emph{Sensors}, 2022. doi: \href{https://doi.org/10.3390/s22145103}{10.3390/s22145103}.

\bibitem{khan} A. A. Khan, A. A. Laghari, M. Shafiq, O. Cheikhrouhou, W. Alhakami, H. Hamam, and Z. A. Shaikh. Healthcare ledger management: A blockchain and machine learning-enabled novel and secure architecture for medical industry. \emph{Human-centric Computing and Information Sciences}, 12, 2022. doi: \href{https://doi.org/10.22967/HCIS.2022.12.055}{10.22967/HCIS.2022.12.055}.

\bibitem{fernando} C. D. Fernando, P. T. Weerasinghe, and C. K. Walgampaya. Heart disease risk identification using machine learning techniques for a highly imbalanced dataset: a comparative study. \emph{KDU Journal of Multidisciplinary Studies}, 4(2), 2022. doi: \href{https://doi.org/10.4038/kjms.v4i2.50}{10.4038/kjms.v4i2.50}.

\bibitem{ishaq} Abid Ishaq, Saima Sadiq, Muhammad Umer, Saleem Ullah, Seyedali Mirjalili, Vaibhav Rupapara, and Michele Nappi. Improving the prediction of heart failure patients' survival using SMOTE and effective data mining techniques. \emph{IEEE Access}, 9:39707--39716, 2021. doi: \href{https://doi.org/10.1109/ACCESS.2021.3064084}{10.1109/ACCESS.2021.3064084}.

\bibitem{singh2015} Arpit Singh and Anuradha Purohit. A survey on methods for solving data imbalance problem for classification. \emph{International Journal of Computer Applications}, 127(15):37--41, 2015. doi: \href{https://doi.org/10.5120/ijca2015906677}{10.5120/ijca2015906677}.

\bibitem{batista} Gustavo Batista, Ronaldo Prati, and Maria-Carolina Monard. A study of the behavior of several methods for balancing machine learning training data. \emph{SIGKDD Explorations}, 6:20--29, 2004. doi: \href{https://doi.org/10.1145/1007730.1007735}{10.1145/1007730.1007735}.

\bibitem{ali} Liaqat Ali, Abdur Rahman, Aurangzeb Khan, MengChu Zhou, Adeeb Javeed, and Jude A. Onwuegbuzie. Efficient prediction of cardiovascular disease using machine learning algorithms with Relief and LASSO feature selection techniques. \emph{IEEE Access}, 9:24505--24522, 2021. doi: \href{https://doi.org/10.1109/ACCESS.2021.3053759}{10.1109/ACCESS.2021.3053759}.

\bibitem{li} Jian Ping Li, Amin Ul Haq, Salah Ud Din, Jalaluddin Khan, Asif Khan, and Abdus Saboor. Heart disease identification method using machine learning classification in e-healthcare. \emph{IEEE Access}, 8:107562--107582, 2020. doi: \href{https://doi.org/10.1109/ACCESS.2020.3001149}{10.1109/ACCESS.2020.3001149}.

\bibitem{fitriyani} Norma Latif Fitriyani, Muhammad Syafrudin, Ganjar Alfian, and Jongtae Rhee. HDPM: an effective heart disease prediction model for a clinical decision support system. \emph{IEEE Access}, 8:133034--133050, 2020. doi: \href{https://doi.org/10.1109/ACCESS.2020.3010511}{10.1109/ACCESS.2020.3010511}.

\bibitem{ahmad} Ghulab Nabi Ahmad, Hira Fatima, Shafi Ullah, Abdelaziz Salah Saidi, et al. Efficient medical diagnosis of human heart diseases using machine learning techniques with and without GridSearchCV. \emph{IEEE Access}, 10:80151--80173, 2022. doi: \href{https://doi.org/10.1109/ACCESS.2022.3165792}{10.1109/ACCESS.2022.3165792}.

\bibitem{ramesh} T. R. Ramesh, Umesh Kumar Lilhore, M. Poongodi, Sarita Simaiya, Amandeep Kaur, and Mounir Hamdi. Predictive analysis of heart diseases with machine learning approaches. \emph{Malaysian Journal of Computer Science}, pages 132--148, 2022. doi: \href{https://doi.org/10.22452/mjcs.sp2022no1.10}{10.22452/mjcs.sp2022no1.10}.

\bibitem{qureshi} Muhammad Bilal Qureshi, Jason Glower, Daniel L. Ewert, and Steven C. Koenig. A novel idea to improve cardiac output of mechanical circulatory support devices by optimizing kinetic energy transfer available in forward moving aortic blood flow. \emph{Cardiovascular Engineering and Technology}, 8(2):131--144, 2017. doi: \href{https://doi.org/10.1007/s13239-017-0305-2}{10.1007/s13239-017-0305-2}.

\bibitem{azmi} Javed Azmi, Muhammad Arif, Md Tabrez Nafis, M. Afshar Alam, Safdar Tanweer, and Guojun Wang. A systematic review on machine learning approaches for cardiovascular disease prediction using medical big data. \emph{Medical Engineering \& Physics}, page 103825, 2022. doi: \href{https://doi.org/10.1016/j.medengphy.2022.103825}{10.1016/j.medengphy.2022.103825}.

\bibitem{singh2022} Vijendra Singh, Vijayan K. Asari, and Rajkumar Rajasekaran. A deep neural network for early detection and prediction of chronic kidney disease. \emph{Diagnostics}, 12(1), 2022. doi: \href{https://doi.org/10.3390/diagnostics12010116}{10.3390/diagnostics12010116}.

\bibitem{singhjain} Vijendra Singh and Divya Jain. A hybrid parallel classification model for the diagnosis of chronic kidney disease. \emph{International Journal of Interactive Multimedia and Artificial Intelligence}, 2021. doi: \href{https://doi.org/10.9781/ijimai.2021.10.008}{10.9781/ijimai.2021.10.008}.

\bibitem{jain2019} Divya Jain and Vijendra Singh. A two-phase hybrid approach using feature selection and adaptive SVM for chronic disease classification. \emph{International Journal of Computers and Applications}, 43:1--13, 2019. doi: \href{https://doi.org/10.1080/1206212X.2019.1577534}{10.1080/1206212X.2019.1577534}.

\bibitem{jain2020} Divya Jain and Vijendra Singh. A novel hybrid approach for chronic disease classification. \emph{International Journal of Healthcare Information Systems and Informatics (IJHISI)}, 15(1):1--19, 2020. doi: \href{https://doi.org/10.4018/IJHISI.2020010101}{10.4018/IJHISI.2020010101}.

\bibitem{mohan} Senthilkumar Mohan, Chandrasegar Thirumalai, and Gautam Srivastava. Effective heart disease prediction using hybrid machine learning techniques. \emph{IEEE Access}, 7:81542--81554, 2019. doi: \href{https://doi.org/10.1109/ACCESS.2019.2923707}{10.1109/ACCESS.2019.2923707}.

\bibitem{kavitha} M. Kavitha, G. Gnaneswar, R. Dinesh, Y. Rohith Sai, and R. Sai Suraj. Heart disease prediction using hybrid machine learning model. In \emph{2021 6th International Conference on Inventive Computation Technologies (ICICT)}, pages 1329--1333, 2021. doi: \href{https://doi.org/10.1109/ICICT50816.2021.9358597}{10.1109/ICICT50816.2021.9358597}.

\bibitem{teboul} Alex Teboul. Heart Disease Health Indicators Dataset. 2015. \url{https://www.kaggle.com/datasets/alexteboul/heart-disease-health-indicators-dataset}.

\bibitem{brfss} Centers for Disease Control and Prevention (CDC). Behavioral Risk Factor Surveillance System (BRFSS), 2022. Available at: \url{https://www.kaggle.com/datasets/cdc/behavioral-risk-factor-surveillance-system?select=2015.csv}.

\bibitem{brfsscalc} Centers for Disease Control and Prevention (CDC). Calculated variables in the 2015 data file of the BRFSS, 2015. Available at: \url{https://www.cdc.gov/brfss/annual_data/2015/pdf/2015_Calculated_Variables_Version4_08_10_17-508c.pdf} (Version 4, 08/10/17).

\bibitem{ullah2022} Zahid Ullah and Mona Jamjoom. An intelligent approach for Arabic handwritten letter recognition using convolutional neural network. \emph{PeerJ Computer Science}, 8:e995, 2022. doi: \href{https://doi.org/10.7717/peerj-cs.995}{10.7717/peerj-cs.995}.

\bibitem{ullah2023} Zahid Ullah and Mona Jamjoom. A deep learning for Alzheimer's stages detection using brain images. \emph{Computers, Materials and Continua}, 74, 2022. doi: \href{https://doi.org/10.32604/cmc.2023.032752}{10.32604/cmc.2023.032752}.

\bibitem{bhatt} Chintan M. Bhatt, Parth Patel, Tarang Ghetia, and Pier Luigi Mazzeo. Effective heart disease prediction using machine learning techniques. \emph{Algorithms}, 16(2), 2023. doi: \href{https://doi.org/10.3390/a16020088}{10.3390/a16020088}.

\bibitem{chang} Wenbing Chang, Yinglai Liu, Xueyi Wu, Yiyong Xiao, Shenghan Zhou, and Wen Cao. A new hybrid XGBSVM model: Application for hypertensive heart disease. \emph{IEEE Access}, 7:175248--175258, 2019. doi: \href{https://doi.org/10.1109/ACCESS.2019.2957367}{10.1109/ACCESS.2019.2957367}.

\bibitem{almazroi} Abdulwahab Ali Almazroi, Eman A. Aldhahri, Saba Bashir, and Sufyan Ashfaq. A clinical decision support system for heart disease prediction using deep learning. \emph{IEEE Access}, 11:61646--61659, 2023. doi: \href{https://doi.org/10.1109/ACCESS.2023.3285247}{10.1109/ACCESS.2023.3285247}.

\bibitem{amarbayasgalan} Tsatsral Amarbayasgalan, Van-Huy Pham, Nipon Theera-Umpon, Yongjun Piao, and Keun Ho Ryu. An efficient prediction method for coronary heart disease risk based on two deep neural networks trained on well-ordered training datasets. \emph{IEEE Access}, 9:135210--135223, 2021. doi: \href{https://doi.org/10.1109/ACCESS.2021.3116974}{10.1109/ACCESS.2021.3116974}.

\bibitem{kaufman} Shachar Kaufman, Saharon Rosset, Claudia Perlich, and Ori Stitelman. Leakage in data mining: Formulation, detection, and avoidance. \emph{ACM Transactions on Knowledge Discovery from Data}, 6(4):15:1--15:21, 2012. doi: \href{https://doi.org/10.1145/2382577.2382579}{10.1145/2382577.2382579}.

\bibitem{kapoor} Sayash Kapoor and Arvind Narayanan. Leakage and the reproducibility crisis in machine-learning-based science. \emph{Patterns}, 4(9):100804, 2023. doi: \href{https://doi.org/10.1016/j.patter.2023.100804}{10.1016/j.patter.2023.100804}.

\bibitem{santos} Miriam Seoane Santos, Jastin Pompeu Soares, Pedro Henriques Abreu, H\'{e}lder Ara\'{u}jo, and Jo\~{a}o Santos. Cross-validation for imbalanced datasets: avoiding overoptimistic and overfitting approaches. \emph{IEEE Computational Intelligence Magazine}, 13(4):59--76, 2018.

\bibitem{deng} Lang Deng, Kongjie Lu, and Huanhuan Hu. An interpretable LightGBM model for predicting coronary heart disease: Enhancing clinical decision-making with machine learning. \emph{PLOS ONE}, 20(9):e0330377, 2025. doi: \href{https://doi.org/10.1371/journal.pone.0330377}{10.1371/journal.pone.0330377}.

\bibitem{hasnat} Md Abrar Hasnat, Md Jobayer, Md. Mehedi Hasan Shawon, and Md. Golam Rabiul Alam. Interpretable heart disease prediction via a weighted ensemble model: A large-scale study with SHAP and surrogate decision trees. \emph{arXiv preprint arXiv:2511.01947}, 2025.

\bibitem{detrano} Andras Janosi, William Steinbrunn, Matthias Pfisterer, and Robert Detrano. Heart Disease [Dataset]. UCI Machine Learning Repository, 1989. doi: \href{https://doi.org/10.24432/C52P4X}{10.24432/C52P4X}.

\bibitem{pmlb} Randal S. Olson, William La Cava, Patryk Orzechowski, Ryan J. Urbanowicz, and Jason H. Moore. PMLB: a large benchmark suite for machine learning evaluation and comparison. \emph{BioData Mining}, 10:36, 2017. doi: \href{https://doi.org/10.1186/s13040-017-0154-4}{10.1186/s13040-017-0154-4}.

\end{thebibliography}
